\documentclass{article}
\usepackage{enumitem}
\usepackage{amsmath}
\begin{document}

\title{Chapter 36: Ecosystem}
\date{}
\maketitle

\begin{enumerate}[label=Q\arabic*.]

\item The 10\% law of energy transfer in ecosystems (Lindeman's law) states that:
\\(A) 10\% of sunlight is fixed by producers
\\(B) Only about 10\% of energy is transferred from one trophic level to the next; 90\% is lost as heat
\\(C) 10\% of organisms in each trophic level survive
\\(D) Energy transfer efficiency is always exactly 10\%

\item Net Primary Productivity (NPP) is defined as:
\\(A) Total energy fixed by producers through photosynthesis (GPP)
\\(B) GPP minus energy lost to plant respiration: NPP = GPP $-$ R
\\(C) Energy available to herbivores at the second trophic level
\\(D) Total biomass produced per year in an ecosystem

\item Which of the following correctly describes the carbon cycle?
\\(A) Carbon moves from atmosphere (CO$_2$) $\rightarrow$ plants (photosynthesis) $\rightarrow$ animals (feeding) $\rightarrow$ atmosphere (respiration, decomposition)
\\(B) Carbon moves only between plants and animals
\\(C) Carbon is entirely recycled through the hydrological cycle
\\(D) Carbon only enters the cycle through combustion of fossil fuels

\item Decomposers (detritivores and saprotrophs) play a crucial role in ecosystems by:
\\(A) Fixing atmospheric nitrogen
\\(B) Releasing nutrients from dead organic matter back into the environment (nutrient cycling)
\\(C) Producing oxygen through photosynthesis
\\(D) Providing food for primary consumers

\item Assertion: The pyramid of energy is always upright and can never be inverted.\\
Reason: Energy is lost at each trophic level as heat; energy at lower trophic levels always exceeds that at higher levels.
\\(A) Both true, Reason is correct explanation
\\(B) Both true, Reason is not correct explanation
\\(C) Assertion true, Reason false
\\(D) Assertion false, Reason true

\item The standing crop refers to:
\\(A) The productivity (energy fixed per unit time)
\\(B) The amount of living biomass present at a given time in a trophic level
\\(C) The crop yield per hectare
\\(D) The number of species in a community

\item Primary succession occurs on:
\\(A) Bare substrate devoid of soil and organisms (e.g., rock, sand dune, lava flow)
\\(B) Areas where existing community was destroyed but soil remains
\\(C) Agricultural land left fallow
\\(D) Burned forest areas

\item Pioneer species in primary succession on bare rock are typically:
\\(A) Trees
\\(B) Lichens (then mosses, then herbs, then shrubs, then trees)
\\(C) Grasses
\\(D) Ferns

\item Eutrophication of freshwater bodies refers to:
\\(A) Introduction of invasive fish species
\\(B) Excessive nutrient (especially N and P) enrichment causing algal blooms, oxygen depletion, and loss of biodiversity
\\(C) Acidification due to acid rain
\\(D) Increase in water turbidity due to silt

\item The phosphorus cycle differs from the nitrogen and carbon cycles in that:
\\(A) Phosphorus does not cycle at all
\\(B) Phosphorus has no atmospheric gaseous phase; it cycles through rocks, soil, water, and organisms only
\\(C) Phosphorus cycles faster than nitrogen
\\(D) Phosphorus is not essential for organisms

\item In a food web, an omnivore occupies:
\\(A) Only the first trophic level
\\(B) Multiple trophic levels simultaneously depending on what it eats
\\(C) Only the top trophic level
\\(D) Only the second trophic level

\item Biological magnification (biomagnification) of toxins like DDT occurs because:
\\(A) DDT is degraded at lower trophic levels
\\(B) DDT accumulates in fat tissue and its concentration increases at each successive trophic level
\\(C) DDT is absorbed only by top predators
\\(D) DDT inhibits enzyme activity at lower trophic levels only

\item Gross Primary Productivity (GPP) of an ecosystem depends primarily on:
\\(A) Number of animal species
\\(B) Photosynthetic capacity of producers, light availability, water, and nutrient supply
\\(C) Rainfall only
\\(D) Soil mineral content only

\item Match the following types of ecological pyramids with their characteristics:\\
1. Pyramid of numbers $\rightarrow$ (a) Can be inverted (e.g., parasitic food chain)\\
2. Pyramid of biomass $\rightarrow$ (b) Can be inverted (e.g., ocean -- phytoplankton vs zooplankton)\\
3. Pyramid of energy $\rightarrow$ (c) Always upright; never inverted\\
4. Inverted pyramid $\rightarrow$ (d) Seen in aquatic ecosystems for biomass
\\(A) 1-a, 2-b, 3-c, 4-d
\\(B) 1-b, 2-a, 3-d, 4-c
\\(C) 1-c, 2-d, 3-a, 4-b
\\(D) 1-d, 2-c, 3-b, 4-a

\item The concept of an ``ecological niche'' encompasses:
\\(A) Only the physical habitat of an organism
\\(B) The total role of an organism in its community including its habitat, food, interactions, and functional position
\\(C) Only the trophic level of an organism
\\(D) The geographic range of a species

\end{enumerate}
\end{document}
